% =====================================================================
%  รายงานการทดสอบ — ตารางสถานการณ์ทดสอบ / Test Case
%  Compile: XeLaTeX  (Overleaf: Menu > Settings > Compiler > XeLaTeX)
%
%  โฟลเดอร์นี้แยกจากเล่มคู่มือ ใช้เก็บ test case ตามรูปแบบตารางในเล่มรายงาน
%  วิธีใช้: กรอกข้อมูลส่วนที่ 1 แล้วเพิ่มไฟล์ใน scenarios/ แล้ว \input ในส่วนที่ 3
% =====================================================================
\documentclass[11pt,a4paper]{article}
\usepackage{testreport}

% ---------------------------------------------------------------------
% 1) ข้อมูลเล่ม
% ---------------------------------------------------------------------
\projectCode{CS 2568/01}
\titleTH{ระบบดูแลและคัดกรองสุขภาพจิตนิสิต (NUCARE)}
\titleEN{NUCARE — Student Mental Health Care and Screening System}
\doctype{รายงานการทดสอบ}

\systemUnderTest{ระบบเว็บแอปพลิเคชัน ``NUCARE''}
\systemVersion{1.0}
\testerName{นายสมชาย ใจดี}
\testDateTH{วันที่ 22 ตุลาคม พ.ศ.\ 2568}

\firstStudentID{653450001-1}
\firstStudentPrefixTH{นาย}
\firstStudentNameTH{สมชาย ใจดี}
\secondStudentID{653450002-9}
\secondStudentPrefixTH{นางสาว}
\secondStudentNameTH{สมหญิง รักเรียน}

\advisorTH{ผศ.ดร.ประสิทธิ์ สร้างสรรค์}
\courseCode{CP354 772}
\courseName{โครงงานคอมพิวเตอร์ 2}
\semester{2}
\academicYearTH{2568}
\submitMonthTH{ตุลาคม}
\submitYearTH{2568}

\begin{document}

% ---------------------------------------------------------------------
% 2) ปก (ลบ \maketestcover ถ้าไม่ต้องการปก)
% ---------------------------------------------------------------------
\maketestcover

\section*{สรุปการทดสอบ}
เอกสารฉบับนี้รวบรวมสถานการณ์ทดสอบ (Test Scenario) และกรณีทดสอบ (Test Case)
ของระบบเว็บแอปพลิเคชัน ``NUCARE'' เพื่อบันทึกขั้นตอน ผลลัพธ์ที่คาดหวัง ผลลัพธ์จริง และผลการทดสอบ
ในรูปแบบตารางมาตรฐานของเล่มรายงานโครงงาน

\begin{center}
\begin{tabularx}{0.92\linewidth}{|c|Y|c|c|}
\hline
\rowcolor{tcHeader}
\textbf{รหัส} & \textbf{ชื่อสถานการณ์ทดสอบ} & \textbf{จำนวนกรณี} & \textbf{ผลรวม} \\
\hline
TC01 & การเข้าสู่ระบบของผู้ใช้งานผ่านหน้าเว็บแอปพลิเคชัน & 2 & Pass \\
\hline
TC02 & การทำแบบประเมินสุขภาพจิตของนิสิตผ่านหน้าเว็บแอปพลิเคชัน & 1 & Pass \\
\hline
\end{tabularx}
\end{center}

\clearpage

% ---------------------------------------------------------------------
% 3) ตารางสถานการณ์ทดสอบ — เพิ่ม \input{scenarios/...} ตามจำนวนเคส
% ---------------------------------------------------------------------
\section*{รายละเอียดสถานการณ์ทดสอบ}
% TC01 — การเข้าสู่ระบบ (รูปแบบตารางตามเล่มอ้างอิง รวมหน้าต่อ)
\begin{testscenario}{1}{1}
  \tcmeta{TC01}{CS 2568/01}{ระบบเว็บแอปพลิเคชัน ``NUCARE''}{นายสมชาย ใจดี}{วันที่ 22 ตุลาคม พ.ศ.\ 2568}{1.0}
  \tcname{การเข้าสู่ระบบของผู้ใช้งานผ่านหน้าเว็บแอปพลิเคชัน}
  \tcpre{ผู้ใช้งานมีบัญชีในระบบแล้ว และสามารถเปิดหน้าเข้าสู่ระบบผ่านเว็บเบราว์เซอร์ได้}
  \tcdesc{ทดสอบการทำงานของระบบเข้าสู่ระบบ โดยกรอกข้อมูลในแต่ละกรณีให้ครอบคลุมทั้งข้อมูลถูกต้องและข้อมูลไม่ถูกต้อง
    เพื่อยืนยันว่าระบบสามารถตรวจสอบสิทธิ์ แสดงข้อความแจ้งเตือน หรือนำผู้ใช้ไปยังแดชบอร์ดได้ถูกต้องตามเงื่อนไขที่ออกแบบไว้}

  \tccase{1}{ตรวจสอบการเข้าสู่ระบบเมื่อผู้ใช้กรอกข้อมูลถูกต้องครบถ้วน}{%
    \begin{tcsteps}
      \item เปิดหน้า ``เข้าสู่ระบบ'' ของระบบ
      \item กรอกชื่อผู้ใช้ = student01
      \item กรอกรหัสผ่าน = Password123!
      \item กดปุ่ม ``เข้าสู่ระบบ''
    \end{tcsteps}}{%
    ระบบตรวจสอบสิทธิ์สำเร็จ และนำผู้ใช้ไปยังแดชบอร์ดตามบทบาท}{%
    ระบบตรวจสอบสิทธิ์สำเร็จ และนำผู้ใช้ไปยังแดชบอร์ดตามบทบาท}{Pass}

  \tccase{2}{ตรวจสอบเมื่อผู้ใช้กรอกชื่อผู้ใช้ว่าง หรือมีความยาวน้อยกว่า 4 ตัวอักษร}{%
    \begin{tcsteps}
      \item เปิดหน้า ``เข้าสู่ระบบ'' ของระบบ
      \item กรอกชื่อผู้ใช้ = ว่าง หรือ stu
      \item กรอกรหัสผ่าน = Password123!
      \item กดปุ่ม ``เข้าสู่ระบบ''
    \end{tcsteps}}{%
    ใต้ช่องชื่อผู้ใช้แสดงข้อความ ``ชื่อผู้ใช้ต้องมีอย่างน้อย 4 ตัวอักษร''}{%
    ใต้ช่องชื่อผู้ใช้แสดงข้อความ ``ชื่อผู้ใช้ต้องมีอย่างน้อย 4 ตัวอักษร''}{Pass}

  \tccasebreak
  \tccasecont{2 (ต่อ)}{%
    \begin{tcsteps}
      \setcounter{enumi}{4}
      \item กรอกยืนยันว่ารหัสผ่านยังเป็น Password123!
      \item กดปุ่ม ``เข้าสู่ระบบ'' อีกครั้งหลังแก้ชื่อผู้ใช้สั้น
    \end{tcsteps}}

  \tccase{3}{ตรวจสอบเมื่อผู้ใช้กรอกรหัสผ่านว่าง หรือรหัสผ่านไม่ถูกต้อง}{%
    \begin{tcsteps}
      \item เปิดหน้า ``เข้าสู่ระบบ'' ของระบบ
      \item กรอกชื่อผู้ใช้ = student01
      \item กรอกรหัสผ่าน = ว่าง หรือ wrong-password
      \item กดปุ่ม ``เข้าสู่ระบบ''
    \end{tcsteps}}{%
    ระบบปฏิเสธการเข้าสู่ระบบ และแสดงข้อความแจ้งว่าชื่อผู้ใช้หรือรหัสผ่านไม่ถูกต้อง}{%
    ระบบปฏิเสธการเข้าสู่ระบบ และแสดงข้อความแจ้งว่าชื่อผู้ใช้หรือรหัสผ่านไม่ถูกต้อง}{Pass}

  \tccase{4}{ตรวจสอบเมื่อผู้ใช้กรอกรหัสผ่านมีความยาวน้อยกว่า 8 ตัวอักษร}{%
    \begin{tcsteps}
      \item เปิดหน้า ``เข้าสู่ระบบ'' ของระบบ
      \item กรอกชื่อผู้ใช้ = student01
      \item กรอกรหัสผ่าน = Pass12
      \item กดปุ่ม ``เข้าสู่ระบบ''
    \end{tcsteps}}{%
    ใต้ช่องรหัสผ่านแสดงข้อความ ``รหัสผ่านต้องมีอย่างน้อย 8 ตัวอักษร''}{%
    ใต้ช่องรหัสผ่านแสดงข้อความ ``รหัสผ่านต้องมีอย่างน้อย 8 ตัวอักษร''}{Pass}
\end{testscenario}

% TC02 — แบบประเมินสุขภาพจิตของนิสิต
\begin{testscenario}{2}{2}
  \tcmeta{TC02}{CS 2568/01}{ระบบเว็บแอปพลิเคชัน ``NUCARE''}{นายสมชาย ใจดี}{วันที่ 22 ตุลาคม พ.ศ.\ 2568}{1.0}
  \tcname{การทำแบบประเมินสุขภาพจิตของนิสิตผ่านหน้าเว็บแอปพลิเคชัน}
  \tcpre{ผู้ใช้งานเข้าสู่ระบบด้วยบทบาทนิสิตแล้ว และมีแบบประเมินที่เปิดให้ทำอยู่ในระบบ}
  \tcdesc{ทดสอบการส่งแบบประเมินสุขภาพจิตเมื่อนิสิตอบครบทุกข้อ
    เพื่อยืนยันว่าระบบบันทึกคำตอบ คำนวณผล และแสดงสรุปผลได้ถูกต้อง}

  \tccase{1}{ตรวจสอบการส่งแบบประเมินเมื่อนิสิตอบครบทุกข้อ}{%
    \begin{tcsteps}
      \item เข้าเมนู ``แบบประเมิน'' จากแดชบอร์ดนิสิต
      \item เลือกแบบประเมินที่ต้องการทำ
      \item ตอบคำถามทุกข้อตามตัวเลือกที่กำหนด
      \item กดปุ่ม ``ส่งคำตอบ''
    \end{tcsteps}}{%
    ระบบบันทึกคำตอบสำเร็จ คำนวณคะแนน/ระดับผลลัพธ์ และแสดงหน้าสรุปผลให้ผู้ใช้}{%
    ระบบบันทึกคำตอบสำเร็จ คำนวณคะแนน/ระดับผลลัพธ์ และแสดงหน้าสรุปผลให้ผู้ใช้}{Pass}
\end{testscenario}


% แม่แบบว่าง: คัดลอก scenarios/_template.tex เป็นไฟล์ใหม่ แล้ว \input ที่นี่

\end{document}
